% =====================================================================
\chapter[可证伪的研究纲领]{可证伪的研究纲领}\label{ch:program}
% =====================================================================

本文的全部主张若不转化为可检验的形式，就只是一组解释框架。本章给出这一转化，其做法
遵循研究纲领方法论中"硬核与保护带"的区分\cite{lakatos1970}：不可承重的主张应被移入
保护带，而不是加固为硬核。

\section{方法论：从解释框架到可检验假设}

\begin{principle}[降格程序]\label{pr:demote}
对每一条主张，执行下述操作：写出它的\emph{可观察后果}；写出什么观察会\emph{否定}
它；写出获得该观察的\emph{最小实验}；写出该实验的\emph{对照条件}。凡是无法完成这
四步的主张，应当被明确标注为解释框架而非经验断言。
\end{principle}

按原则 \ref{pr:demote}，本文的主张被分为三类：

\begin{itemize}
\item \term{已证明的}：命题 \ref{prop:oblivious}、\ref{prop:overconstrain}、
      \ref{prop:stopping}、\ref{prop:leakpair}、\ref{prop:soundness}、
      \ref{prop:active}、\ref{prop:asym-struct}、\ref{prop:operad-compose}、
      \ref{prop:transport}、\ref{prop:observability}、\ref{prop:exponent}；定理
      \ref{th:integration}、\ref{th:six}、\ref{th:centralise}、
      \ref{th:condorcet}；命题 \ref{th:correlated}；推论
      \ref{cor:threshold}、\ref{cor:buildadopt}、\ref{cor:diversity}、
      \ref{cor:three-failures}。
\item \term{已论证但依赖明确假设的}：原则 \ref{th:compression}、
      \ref{th:necessity}、\ref{th:separation}、\ref{th:design}；命题
      \ref{prop:no-godel}、\ref{prop:three-grounds}。
\item \term{作为可证伪假设提出的}：下述 $H_1$--$H_{14}$。此外，注
      \ref{rmk:fibration-status} 的结构提案（把层族提升为严格纤维化）与
      注 \ref{rmk:rg} 的重整化类比均未被计入定理，其地位是提案与启发式。
\end{itemize}

\section{十四条可证伪假设}

\begin{hypothesis}[结构收敛高于随机基线]\label{hyp:h1}
在控制系统规模、领域、平台、编程语言与框架之后，来自不同领域的大型软件系统在其
抽象骨架（定义 \ref{def:skeleton}）上呈现的结构相似性显著高于随机基线。
\emph{操作化}：以第 \ref{sec:baseline} 节的五步设计构造架构图集合与骨架映射，
以第三节的结构指标为特征向量，检验跨领域配对的相似性分布是否随机优于此零模型
（如度序列保持的图重连）。
\emph{证伪条件}：相似性分布与零模型无显著差异。
\emph{风险}：骨架映射的选择自由度；需预注册映射与指标。
\end{hypothesis}

\begin{hypothesis}[收敛强度关于领域绑定度单调]\label{hyp:h2}
猜想 \ref{conj:strength} 成立：结构相似性关于领域绑定度
$\mathrm{Dom}$（定义 \ref{def:dom}）单调递减。
\emph{操作化}：对每个层级标注 $\mathrm{Dom}$（由该层语义所依赖的领域公理占比估计），
回归相似性对 $\mathrm{Dom}$。
\emph{证伪条件}：单调性不成立或系数不显著。
\end{hypothesis}

\begin{hypothesis}[相变阈值由 $n \times m$ 决定]\label{hyp:h3}
系统是否建立统一基座，由活跃实体数与共享资源数的乘积（以及重复问题频率与演化频率）
预测得比代码行数更好。
\emph{操作化}：对一组开源项目，人工或半自动标注（是否出现六职能中的自建实现），
以 $n \times m$ 的代理量（并发实体数量级、共享资源种类数）与代码行数分别做逻辑回归，
比较解释力。
\emph{证伪条件}：代码行数的解释力不劣于 $n \times m$。
\emph{风险}：$n$ 与 $m$ 的自动估计困难；可先用并发实体与互斥量/连接池/缓存等的
数量作为代理。
\end{hypothesis}

\begin{hypothesis}[读集的标度优于传统规模指标]\label{hyp:h4}
变更读集 $\mathbb{E}[\rho]$（定义 \ref{def:readset}）关于系统规模的\emph{标度指数}
$\alpha$（命题 \ref{prop:exponent}）与缺陷率、变更失败率的相关性，优于模块数、
代码行数、圈复杂度等传统规模指标，也优于读集的当前水平
$\mathbb{E}[\rho]$ 本身。
\emph{操作化}：在若干有完整提交历史与缺陷追踪的开源项目上，按时间窗口同时估计读集
的水平与标度，分别与缺陷密度做面板回归，比较两者的解释力。
\emph{证伪条件}：标度不提供超越水平与既有指标的预测力。
\end{hypothesis}

\begin{hypothesis}[分层模型加主动探测优于穷尽阅读]\label{hyp:h5}
在单位成本（token 预算或人时）约束下，采用分层认知模型加主动探测的理解过程，
其任务成功率（以定义 \ref{def:understand} 的 (U1)--(U5) 为评分标准）优于穷尽式阅读。
\emph{操作化}：给定同一陌生代码库与同一组任务（预测影响面、定位缺陷、添加功能），
比较三种策略：顺序阅读、仅层次摘要、层次摘要加探测循环。测量任务成功率与单位成本
成功率。
\emph{证伪条件}：主动探测策略不显著更优，或优势完全来自探测次数而非策略结构。
\end{hypothesis}

\begin{hypothesis}[泄漏对可被自动发现]\label{hyp:h6}
命题 \ref{prop:leakpair} 的泄漏对可以被自动化差分测试发现，且发现率与模块边界质量
负相关。
\emph{操作化}：对若干库，生成"接口相同、内部状态不同"的执行对，检测外部行为分歧；
与人工标注的边界质量评分比较。
\emph{证伪条件}：自动生成的执行对无法有效触发泄漏，或发现率与边界质量无关。
\end{hypothesis}

\begin{hypothesis}[规格质量是当前的约束瓶颈]\label{hyp:h7}
在当前技术水平下，限制"意图到可验证执行"的主要瓶颈是规格质量，而非代码生成能力或
证明搜索能力。
\emph{操作化}：以既有基准证据（规格生成任务、代码正确率与规格质量分离评测、
证明成功率）为初始证据，并在新的受控基准上比较三个环节的改进对端到端成功率的
边际贡献。
\emph{初始证据}：在 189 个 Lean 任务上的分项评测报告最佳模型的代码正确率
72.6\%、规格健全性与完备性 52.3\%、证明成功率 4.9\%（单次试验）\cite{verina2026}；
245 个操作系统内核规格生成任务的评测显示模型在长上下文与规格完备性上仍显不足
\cite{osvbench2026}；以可信测试实例化规格的评测指出证明难度会干扰对规格质量的评价
\cite{coins2026}；可追溯精化方法报告在规格生成子任务上的准确率提升
\cite{verispecgen2026}；变异驱动方法报告精确率与召回率\cite{specsyn2026}。
这些是特定基准上的实证结果，不是通用可靠性定理，但它们支持"三个维度应分别评估"
这一判断。
\emph{证伪条件}：在控制其他环节能力后，规格质量的边际贡献不显著高于其他环节。
\end{hypothesis}

\begin{hypothesis}[共识可靠性受相关错误模式支配]\label{hyp:h8}
在意图审定任务上，增加验证者数量对可靠性的提升受相关错误模式 $q$ 的上限约束
（命题 \ref{th:correlated}）；引入异质验证者（不同背景的人、不同来源的模型、
对抗性角色）对可靠性的提升优于单纯增加同质验证者。
\emph{操作化}：构造意图审定任务，配置若干组验证者（同质/异质），测量多数判断的
错误率随组大小与异质性的变化，并与式 \eqref{eq:chernoff} 与 \eqref{eq:correlated}
的预测比较。
\emph{证伪条件}：同质验证者的错误率随 $n$ 下降的速率与异质验证者无显著差异
（即 $q$ 实际上很小）。
\end{hypothesis}

\begin{hypothesis}[制品绑定缺口是可测量的风险]\label{hyp:h9}
"被验证的工件"与"实际执行的工件"之间的绑定缺口（式 \eqref{eq:binding}）在真实
部署中是可测量且可显著降低的风险，且其降低带来可观测的可靠性收益。
\emph{操作化}：在若干部署流水线上引入内容寻址、可复现构建与部署时校验，测量
构建不确定性比例、部署漂移事件率与相关事故率的变化。
\emph{证伪条件}：缺口在实践中不可测量，或修复成本远超收益。
\end{hypothesis}

\begin{hypothesis}[六职能是粗粒化下的存活集]\label{hyp:h10}
在定义 \ref{def:coarsegrain} 的粗粒化序列下，存活的结构恰为 $\mathcal{K}_6$ 的槽位
及其关系型（猜想 \ref{conj:relevant}）。
\emph{操作化}：按原则 \ref{pr:coarsegrain-impl} 对若干仓库实施机械粗粒化，记录存活集
随粒度的变化。
\emph{证伪条件}：存活集包含六职能之外的结构，或六职能中的某些在中等粒度上即被冲掉。
\emph{风险}：粗粒化的阈值选择会引入人为参数；应对阈值做敏感性分析。
\end{hypothesis}

\begin{hypothesis}[激励结构降低相关错误模式]\label{hyp:h11}
在意图审定任务上，采用适当评分规则（原则 \ref{pr:scoring}）相比仅增加验证者数量，
能更显著地降低多数判断的错误率，并可测得错误相关性的下降。
\emph{操作化}：对同一组规格审定任务，比较三种配置：少量无激励验证者、大量无激励
验证者、少量带评分规则验证者；测量错误率与验证者间误差相关系数。
\emph{证伪条件}：评分规则配置的错误率与误差相关系数均不优于大量无激励配置。
\end{hypothesis}

\begin{hypothesis}[率失真斜率预测重构收益]\label{hyp:h12}
率失真曲线的斜率 $\partial L_T / \partial \DL$（注 \ref{rmk:rd-refactor}）能预测
哪些重构在实际中降低缺陷率或返工率。
\emph{操作化}：对一组已完成的重构，事前估计其沿曲线的位移方向与斜率，事后以缺陷率
与变更失败率的变化作为结果变量做回归。
\emph{证伪条件}：斜率不提供超越"改动规模"等朴素变量的预测力。
\end{hypothesis}

\begin{hypothesis}[单位代码量的证明成本轨迹]\label{hyp:h13}
第 \ref{ch:objections} 章并不主张形式化验证会\emph{自动}变便宜，而是要求一个可测量的
成本指标。该假设是：比值
\[
\kappa_{\text{proof}} \;=\; \frac{\text{确立一项已陈述保证所需的工作量}}{\text{承载该保证的代码规模}}
\]
处在一条可测量的轨迹上，并且决定机械化验证会变成常规工程、还是仍被局限在安全关键
工作中的，是这条轨迹——而不是它的当前取值。
\emph{初始证据}：seL4 的功能正确性证明据报告约用十一个人年，对应约 8{,}500 行 C 代码，
即约每千行 $1.3$ 人年，且\emph{不含}对可复用工具与框架的投入 \cite{klein2009}；把该
投入计入的数字要大好几倍。这两个数字之间的差距值得点明，因为可复用工具所占的那一份，
恰恰是该假设成立时应当下降的部分。
\emph{操作化}：对固定的性质类与固定的证明器，按项目、按年度记录
$\kappa_{\text{proof}}$，把可复用基础设施所占份额与单项目所占份额分开；在观测窗口上
拟合这条轨迹。
\emph{证伪条件}：测不出趋势，或者趋势完全归因于可复用基础设施份额、而单项目份额
持平。由于本文拒绝预测它朝哪个方向走，两种结果都是有信息量的。
\emph{风险}：分母是有争议的——代码行数、规格行数还是证明行数——这一选择必须预注册，
否则轨迹在不同项目之间不可比。
\emph{解释区间}：为读取这条轨迹而按惯例提出的阈值，例如"常规"的 $0.1$ 与"专家"的
$0.5$--$1$，是约定的判定区间而非测量量，本文在此照此标注它们。
\end{hypothesis}

\begin{hypothesis}[智能体式生成能否降低读集的标度指数]\label{hyp:h14}
命题 \ref{prop:exponent} 使读集的标度指数 $\alpha$ 成为决定一个系统的治理成本随其
规模线性增长还是二次增长的量。上述每一个假设涉及的都是由人、或由人指挥的工具所构建
的系统。这一条直接问出 AI 时代的问题：一个智能体式生成器在被允许重构的条件下，能否
\emph{发现} $\alpha$ 低于人类基线的架构，还是它只是在所继承的架构内部移动常数？这一
区分不是修饰性的——它是"加速工作的自动化"与"改变规模上仍然可能之事的自动化"之间的
区别。
\emph{操作化}：让一个智能体组与一个人工对照组从同一代码库、同一功能待办清单、同样的
重构自由度出发；在同一窗口内跟踪 $\mathbb{E}[\rho]$ 与拟合出的 $\alpha$ 关于系统规模
的变化；比较两条轨迹。
\emph{证伪条件}：$\alpha$ 在两组之间无显著差异，或者表面差异完全由所写代码的数量而非
其结构解释。
\emph{风险}：$\alpha$ 是通过定义 \ref{def:readset} 的读集代理估计的，而该代理的有效性
本身就是假设 $H_4$；$H_4$ 的代理若不成立，会使本假设变得\emph{不可证伪}而不是
\emph{为假}，两者不可混为一谈。
\end{hypothesis}

\begin{table}[htbp]
\centering
\footnotesize
\caption{假设与实验设计汇总}\label{tab:hypotheses}
\begin{tabular}{L{0.9cm}L{4.0cm}L{4.6cm}L{4.6cm}}
\toprule
\textbf{编号} & \textbf{假设} & \textbf{主要测量} & \textbf{主要风险} \\
\midrule
$H_1$ & 结构收敛高于随机基线 & 骨架相似性 vs 零模型 & 骨架映射选择自由度过大 \\
$H_2$ & 收敛强度关于领域绑定度单调 & 相似性对 $\mathrm{Dom}$ 的回归 &
  $\mathrm{Dom}$ 难以独立标注 \\
$H_3$ & 阈值由 $n\times m$ 决定 & 逻辑回归解释力比较 & $n$、$m$ 的自动估计 \\
$H_4$ & 读集标度优于传统规模指标 & 读集标度与缺陷密度的面板回归 & 读集代理精度 \\
$H_5$ & 分层加主动探测优于穷尽阅读 & 单位成本任务成功率 & 任务与评分的主观性 \\
$H_6$ & 泄漏对可被自动发现 & 差分测试发现率 & 触发有效执行对的难度 \\
$H_7$ & 规格质量是瓶颈 & 三环节的边际贡献 & 基准的真值与覆盖 \\
$H_8$ & 共识受 $q$ 支配 & 同质/异质验证者错误率 & 异质性的构造与度量 \\
$H_9$ & 制品绑定缺口可测量 & 构建不确定性与漂移事件率 & 真实部署的可得性 \\
$H_{10}$ & 六职能是粗粒化存活集 & 粗粒化序列下的存活集 & 粗粒化阈值的人为性 \\
$H_{11}$ & 激励结构降低相关错误 & 错误率与误差相关系数 & 激励机制的实验伦理与成本 \\
$H_{12}$ & 率失真斜率预测重构收益 & 重构前后的缺陷率回归 & 斜率的可估计性 \\
$H_{13}$ & 单位代码量的证明成本呈可测量轨迹 &
  $\kappa_{\text{proof}}$ 的拟合轨迹，工具份额单列 & 分母有争议；跨项目可比性 \\
$H_{14}$ & 智能体式生成可降低读集标度指数 $\alpha$ &
  拟合 $\alpha$ 关于系统规模的两组轨迹 & 依赖 $H_4$ 的读集代理 \\
\bottomrule
\end{tabular}
\end{table}

\section{三个优先研究方向}

以下三个方向分别连接需求工程与形式化方法、编译正确性与资源分析、程序契约与运行时
安全。它们的共同特征是：都有明确的精确问题、可证伪假说、检验方法与主要风险。

\subsection*{方向一：面向真实意图的规格质量证据链}

\paragraph{精确问题。}
给定自然语言需求集合 $R$、领域模型 $D$ 与生成规格 $S$，如何构建独立评估流程，
分别检查：需求覆盖；禁止行为排除；合法行为保留；环境假设显式性；规格的可满足性与
可实现性？

\paragraph{已有基础。}
需求追溯研究、需求形式化与验证、空洞性检测、以及近期的规格生成与联合评测基准。
关键连接是需求工程的\emph{语义追溯}与形式化方法的\emph{规格与程序反例}。

\paragraph{可证伪假说。}
与单一"证明通过"或单一"参考规格相似度"相比，原子需求追溯加正反向测试加规格强度与
变异检查加独立领域审查，能更早识别规格的\emph{过弱}、\emph{过强}与\emph{意图遗漏}
（即假设 $H_7$ 的具体化）。

\paragraph{检验方法。}
在支付、权限控制与幂等接口等小型基准上，建立由多利益相关者标注的需求原子、允许与
禁止轨迹集以及环境前提；比较四类方案：纯生成、需求分解、追溯测试、变异精化加专家
辅助。测量精确率与召回率、错误类别、漏掉的禁例率、过度约束率、专家一致性与成本。

\paragraph{证伪条件。}
若加入追溯与变异层并未显著降低隐藏负向要求的漏报率，或收益只来自泄露的测试参考
答案，则假说不成立。

\paragraph{主要风险。}
真实意图的"金标准"本身可能存在争议。

\subsection*{方向二：面向实际制品的功能--资源联合翻译证书}

\paragraph{精确问题。}
给定源程序 $P$、功能契约 $\varphi$、资源预算 $B$ 与固定目标平台模型 $H$，能否为
\emph{每次编译}生成证书 $C$，使独立验证器同时检查语义保持与可分析的资源约束？

\paragraph{已有基础。}
经验证编译器的语义保持证明、逐次翻译验证、翻译认证、证明携带代码、自动摊还资源分析
与最坏执行时间分析。实质连接是把目前\emph{分别}验证的"行为关系"与"成本关系"置于
\emph{带成本的行为语义}中。

\paragraph{工作模型。}
程序语义扩展为二元组 $(\tau, c)$，其中 $\tau$ 是可观察迹线、$c$ 是成本；翻译证明
应保证目标迹线可由源语义解释，并在明确成本模型下证明 $c_{\text{target}} \le B$，
或给出相对于源成本的保守关系。

\begin{remark}[一个必须避免的假设]
不能预设优化编译后的成本一定单调下降。优化可能提高平均性能而恶化最坏情况，
因此成本关系必须是双向可比较的保守关系，而非单调性假设。
\end{remark}

\paragraph{检验方法。}
选定可复现的子语言与目标架构，比较三条路线：经验证编译、翻译验证、普通编译加事后
验证。记录证书大小、验证时间、程序类别覆盖、资源上界紧致度、误拒率以及对编译优化
变化的稳健性。

\paragraph{证伪条件。}
若资源证明在真实平台校准后无法保持可靠，或证书生成与验证成本远高于基线且无安全
收益，则应缩小到特定安全关键领域。

\paragraph{主要风险。}
目标处理器微架构与调度模型不够精确。

\subsection*{方向三：不可信组件的组合契约与运行时保障}

\paragraph{精确问题。}
当一个服务由已验证组件、未完全验证的生成组件与外部服务组成时，如何用假设--保证
契约、信息流策略、资源预算与运行时监控，给出可审计的整体安全命题及故障影响上界？

\paragraph{已有基础。}
模块化假设--保证推理、动态系统契约组合、信息流非干扰、已验证内核与运行时验证的
不确定性研究。实质连接是把子组件的形式化契约与运行时对契约前提与保证的检测及强制
执行放到同一个系统论证中。

\paragraph{可证伪假说。}
当所有组件都有显式假设、出入口事件可观测、且违反契约可由独立监控触发隔离或回退时，
能在不证明每个组件内部实现正确的条件下，为指定安全性质给出可组合的条件性保证。

\paragraph{检验方法。}
搭建多组件支付流程：请求幂等服务、身份权限服务、生成式风险决策组件、外部数据库与
网络模拟器。对比三组（无契约、静态契约、契约加监控与回退）在故障注入下测量重复
扣款、跨用户数据泄露、检测延迟、隔离成功率与性能代价。

\paragraph{证伪条件。}
若监控所能观测的事件不足以判断关键性质，或回退机制自身无法保持安全，则架构不能
提供宣称的保证。

\paragraph{主要风险。}
对于未观测的微妙语义性质，运行时保障仅能约束可检测与可强制执行的部分。

\section{六阶段实施路线}

\begin{table}[htbp]
\centering
\small
\caption{研究实施路线}\label{tab:roadmap}
\begin{tabular}{L{1.0cm}L{6.3cm}L{7.0cm}}
\toprule
\textbf{阶段} & \textbf{核心任务} & \textbf{交付物} \\
\midrule
1 & 定义基准与威胁模型：选择支付服务与幂等接口作为贯穿案例；明确资产、角色、
  可观察事件、禁止行为、环境假设与时限 &
  需求原子集、正负轨迹集、环境假设清单（经利益相关者审定） \\
2 & 规格质量评估：比较纯生成、追溯精化、变异与反例驱动方案；检查过弱规格、过强规格、
  空洞成立、不可满足与不可实现 &
  规格质量评估器；每个需求标识对应规格条款、测试与审定记录 \\
3 & 有限语言程序综合：选择明确子语言与判定性可控的逻辑；分离候选生成与证明检查 &
  候选程序、证明义务、反例分类；区分不可满足、未知、超时与不可表达 \\
4 & 制品与资源证明：选择成熟编译或翻译验证路径；指定可复现实验硬件或仿真模型 &
  编译链元数据、语义保持或翻译证书、资源上界与模型边界 \\
5 & 组合与运行时保障：为外部服务与未完全验证模块写假设--保证契约；添加独立监控、
  沙箱、审计与安全回退 &
  系统级契约组合结果、故障注入报告、剩余风险清单 \\
6 & 端到端评价：测量保证覆盖、漏报、误报、开发成本、证明成本、制品验证成本、
  资源预算与失效恢复 &
  一个能证伪研究假说的公开基准、工具原型与实验报告 \\
\bottomrule
\end{tabular}
\end{table}

\section{成功判据}

以下判据评估该路线是否成功。它们被设计为可独立检验的形式：

\begin{enumerate}[label=(S\arabic*)]
\item \term{对象明确}：每一段保证的对象（哪个工件、哪个语义、哪个环境）可以被唯一
      标识。
\item \term{假设真实}：所依赖的环境假设可以被检验，而非被默认。
\item \term{证据独立}：证据的判定不依赖生成者的自我评价。
\item \term{失败类别可区分}：不可满足、未知、超时、不可表达与候选被反例否定必须被
      分别报告，不能合并为一个"失败"。
\item \term{成本可测量}：生成成本、证明成本、验证成本与运行成本可以被分别计量。
\item \term{可证伪}：至少有一条 $H_i$ 的否定会导致该纲领的核心主张被削弱。
\end{enumerate}

\begin{remark}[纲领的自我检验]
本文的一个核心主张是"生成不可信、验证须独立"。该原则同样适用于本文：本文的承重
命题必须可被独立检验，且不得依赖任何未证明猜想。注 \ref{rmk:q1q8} 与
第 \ref{ch:objections} 章的不作主张清单，是这一自查的具体结果。
\end{remark}

\section{最小可行子集与采纳风险}\label{sec:mvs}

上述假设共享一个它们都没有覆盖的失败模式。它们各自会被实验否证或确证，而本框架还可能
干脆没有被采纳，此时根本不会有实验被运行。这一风险值得明确写出，因为它与可证伪性
正交，也因为它有具体的形状：形式对象是有维护成本的——声明必须被写出、保持更新并被
辩护——如果这个成本超过收益，这套实践就会死于被忽视，而不是死于被否证。

\begin{principle}[最小可行子集]\label{pr:mvs}
如果本框架中只有三个要素能够被采纳，它们应当是：
\begin{enumerate}[label=(\alph*)]
\item \term{无视集声明}：每个层级或模块声明它被允许永远不必关心什么（定义
      \ref{def:oblivious}）；
\item \term{泄漏对测试}：对每一条被声称的边界，用差分测试寻找两条接口轨迹相同而外部
      行为不同的执行（命题 \ref{prop:leakpair}、推论 \ref{cor:diagnose}）；
\item \term{验证义务的层级归属}：每个性质被分派给能够表达它的最低层级，这一分派决定
      它是机械可判定的、可监控的，还是需要审定的（定义 \ref{def:obligation-assign}）。
\end{enumerate}
这三者构成本框架中最小的闭环：一个声明、一个对该声明的测试，以及一个消费前两者的
路由决定。第 \ref{sec:operad} 节的契约操作元与第 \ref{sec:consensus} 节的共识界是
精化层——在相应困难存在之处有价值，在不存在之处可以省略。原则 \ref{pr:descent} 与
\ref{pr:two-criteria} 使这个闭环可用：前者说明要看得多深，后者说明声明必须满足什么。
\end{principle}

采纳主张的可测量形式被刻意写得谦逊，而它正是第 \ref{ch:levels} 章已经为其自身经验
内容提出的指标：本应被机械保证的性质却被留给人工审查的比率，以及被声称的边界未能
通过其自身测试的比率。如果在真实项目中采纳上述三个要素并不能让这些比率移动，那么本
框架的问题就不是某个假设为假，而是这副药比病本身更贵。

把采纳风险写成第十五条假设也是站得住的，以采纳率与成本收益比作为它的测量量。本文没有
这样做，理由本身也是本框架纪律的一部分：一条假设的否证应当取决于关于世界的观察，而
未被采纳取决于实践者的决定，让同一个词同时覆盖两者，会模糊第 \ref{ch:intent} 章着力
保持清晰的"验证"与"确认"之分。这一选择被记录在此，而不是被默默做出。
